% s08 Reflection positivity II: the limit, OS3 and the polynomial form
\section{Reflection positivity II: the limit, \texorpdfstring{$\OSth$}{OS3} and the polynomial form}

The lattice statements of the previous section are now carried to the limits $k\to\infty$ and $J\to\infty$. The scheme of proof is the same for the curvature sector and for the Wick tube: uniform local bounds for the modified-box functionals at complex coefficients (\cref{lem:s08:claim4}), Montel's theorem in $k$ at fixed $J$, the two-sided sandwich of \cref{lem:s07:L,lem:s07:U} on the cone as $J\to\infty$ with the identity theorem off it, and a translation for the planes that are not mid-planes. The sector-specific input is the convergence of the cube functionals themselves: the convergence at fixed cubes (\cref{def:s03:conv}; on King's convention \cref{hyp:s03:EC}) in the curvature sector, and on the Wick tube the positive-leg bound of \cref{cor:s11:W2prime} together with \cref{hyp:s03:W1} (\cref{prop:s08:T}). The arbitrary-amplitude form of the Wick sector, which needs \Wtwo, is the optional extension of \cref{thm:s23:OS2W}; the one lemma written for it here, \cref{lem:s08:W1}, is marked as such and is not used by the main theorem.

\subsection{The enlarged family}

\begin{definition}[the enlarged family, its tube and its bounded domain]\label{def:s08:enlarged}
Fix $\nu\in\{1,2,3\}$, a mid-plane $c\in D$ (\cref{def:s07:reflection}), and finite families $\zeta_1,\dots,\zeta_N$ of real sources $f_{il}\in C^\infty_c(\R^3;\Lambda^2)$, $h_{il}\in C^\infty_c(\R^3)$ supported in $\{x_\nu\ge c+\delta\}$ for some $\delta>0$, with complex coefficients $z_{il},w_{il}$. Split each scalar source as $h_{il}=h^+_{il}-h^-_{il}$ with $h^-_{il}\defeq\norm{h_{il}}_\infty\chi_{il}$, where $\chi_{il}\in C^\infty_c$, $0\le\chi_{il}\le1$, $\chi_{il}=1$ on $\supp h_{il}$ and $\supp\chi_{il}\subset\{x_\nu\ge c+\delta/2\}$, and $h^+_{il}\defeq h_{il}+h^-_{il}\ge0$. The \emph{enlarged family} is the family of real sources
\[
\cF\defeq\{\rho^*f_{il}\}\cup\{f_{il}\}\cup\{h^\pm_{il}\circ\rho\}\cup\{h^\pm_{il}\},\qquad\text{coefficient space }\C^m,\ m=m_F+m_h ,
\]
whose scalar members $\tilde h_j\ge0$ are all nonnegative. For a real coefficient vector $\xi\in\R^m$ the total source is $(f_\xi,k_1(\xi)-k_2(\xi))$ with $k_1(\xi)\defeq\sum_j(\xi_j)_+\tilde h_j$, $k_2(\xi)\defeq\sum_j(\xi_j)_-\tilde h_j$, both smooth and nonnegative ($x_\pm\defeq\max(\pm x,0)$). Let $J_0$ be the least $J$ such that every member of $\cF$ is supported at $\ell^\infty$-distance $>1$ from $\R^3\setminus C_J$; for $J\ge J_0$, $n=n(\lmu,c)$ and the boxes $\Omega_{J,k}$ of \eqref{eq:s07:margin} put
\[
\Phi_k(\zeta)\defeq Z^{a,(3.38)}_{k,\Omega_{J,k}}(\zeta),\qquad\zeta\in\C^m ,
\]
an entire function (a finite-dimensional integral with Gaussian decay in $A$ and $e^{-\lambda\eps^3r^4}$ decay in $r$). The \emph{matrix points} $\zeta^*_{ij}\in\C^m$ carry $\bar z_{il}$ on $\rho^*f_{il}$, $z_{jl}$ on $f_{jl}$, $\pm\bar w_{il}$ on $h^\pm_{il}\circ\rho$, $\pm w_{jl}$ on $h^\pm_{jl}$ and $0$ elsewhere; their total source is that of $\theta\zeta_i^*\sqcup\zeta_j$, so $\Phi_k(\zeta^*_{ij})$ is the $(i,j)$ entry of \cref{lem:s07:RPbox}(b). Since reflected and unreflected members have disjoint supports, $\norm{h^+}_\infty\le2\norm h_\infty$ and $\norm{h^-}_\infty=\norm h_\infty$,
\begin{equation}\label{eq:s08:ampbound}
\norm{k_1(\Re\zeta^*_{ij})}_\infty\ \le\ 2\max\big(\amp(\zeta_i),\amp(\zeta_j)\big).
\end{equation}
\begin{enumerate}[label=\textup{(\roman*)},leftmargin=*,nosep]
\item The \emph{Wick tube} of $\cF$ is $T_\cF\defeq\{\zeta\in\C^m:\ 2\norm{k_1(\Re\zeta)}_\infty<\sigma_+/2\}$, with $\sigma_+=s_1$ the amplitude of \Itwo\ (\cref{def:s10:I2}, \eqref{eq:s11:constants}). It is open, convex ($\xi\mapsto\norm{k_1(\xi)}_\infty$ is a supremum of convex functions), connected, invariant under $z\mapsto-z$ on the curvature coordinates, contains $0$, and by \eqref{eq:s08:ampbound} contains every $\zeta^*_{ij}$ as soon as $\amp(\zeta_i)<\sigma_+/8=:\sigma_W$ for every $i$.
\item The \emph{bounded domain}: let $B_1\subset\R^m$ be the open convex hull of the unit $\ell^\infty$-balls about $0$ and about the real parts $\Re\zeta^*_{ij}$; let $\beta_j\defeq2\sup_{\xi\in B_1}|\xi_j|$ for the scalar coordinates (so $\beta_j=2(|\Re w_{il}|+1)$ for the coordinate of $h^\pm_{il}$ or of $h^\pm_{il}\circ\rho$), $B''\defeq\prod_j[0,\beta_j]$ over the scalar coordinates, and $q_v\defeq\sum_jv_j\tilde h_j$ for $v\in\vertex(B'')$, the \emph{vertex sources}. All $q_v$ are pointwise $\le q_\zeta\defeq\sum_{i,l}2(|\Re w_{il}|+1)(h^+_{il}+h^-_{il}+h^+_{il}\circ\rho+h^-_{il}\circ\rho)$. Put $\cT_1\defeq\{\zeta\in\C^m:\Re\zeta\in B_1\}$: open, convex, connected, containing every $\zeta^*_{ij}$ and the unit polydisc about $0$, so that $\cT_1\cap T_\cF$ is a non-empty open set. For $\xi\in B_1$ both $2\xi_+$ and $\xi_+$ lie in $B''$.
\item The \emph{uniqueness set} $U\subset\C^m$ is the set of coefficient vectors whose curvature coefficients are purely imaginary and whose scalar coefficients are positive reals. It is a non-empty, relatively open subset of the totally real subspace $(\mathrm i\R)^{m_F}\times\R^{m_h}$ of maximal real dimension $m$ (it is not open in $\C^m$), hence a set of uniqueness for holomorphic functions on any connected open set meeting it: a holomorphic function vanishing on a relatively open piece of a maximal totally real subspace has all derivatives zero there \cite{Hor90}. $U\cap\cT_1\ne\emptyset$ and $U\cap T_\cF\ne\emptyset$ (small positive scalar coefficients). On $U$ the total source is $(\mathrm if_\zeta,g_\zeta)\in\cK$, so $\Theta(\zeta)\in\cI$; the scheme multipliers are positive reals, equal for the Neumann--Dirichlet and the modified states; and every lattice value is real (charge conjugation, \cref{def:s02:sources}: $\overline{Z^a(\mathrm if,g)}=Z^a(-\mathrm if,g)=Z^a(\mathrm if,g)$).
\end{enumerate}
Every lattice functional, and hence every limit below, depends on a family only through its total sources; the splitting is therefore immaterial at the matrix points and at the Taylor coefficients at $0$.
\end{definition}

\subsection{Uniform local bounds}

The following local bound is used only in the optional extension (\cref{thm:s23:OS2W}); it is stated here because its domain $\cT_1$ is what makes ``finitely many instances of \Wtwo'' literally true there. The fixed-cube limits $Z^a_J\defeq\lim_kZ^a_{k,J}$ exist for all complex coefficients and are entire under \Wtwo(a) alone: it gives convergence at the real points, the vertex argument of \cref{lem:s08:claim4} with \cref{lem:s02:alg}(d) gives local boundedness in $k$, and Montel's theorem with the uniqueness set $\R^m$ does the rest.

\begin{lemma}[uniform local bound; optional extension]\label{lem:s08:W1}
Assume \Wtwo(a) and \Wtwo(b) at the finitely many sources $q_v$, $v\in\vertex(B'')$. Then $\{Z^a_J\}_{J\ge J_0}$ is uniformly bounded on $\cT_1$, $Z^a_J\to Z^a_\infty$ locally uniformly on $\cT_1$, $Z^a_\infty$ is holomorphic on $\cT_1$, and $Z^a_J\le Z^a_\infty$ on $U\cap\cT_1$.
\end{lemma}

\begin{proof}
For real $\xi\in B_1$, \cref{lem:s02:alg}(f) gives $Z^a_J(\xi)\le Z^a_J(2f_\xi,-2k_2(\xi))^{1/2}Z^a_J(0,2k_1(\xi))^{1/2}$. The first factor is nonincreasing in $J$: at fixed $k$ this is the second inequality of \cref{lem:s02:cut} with \cref{lem:s02:alg}(a), and it passes to $k\to\infty$ by \Wtwo(a); \Wtwo(b) is not used there. Hence it is $\le Z^a_{J_0}(2f_\xi,-2k_2(\xi))$, and $\xi\mapsto\ln Z^a_{J_0}(2f_\xi,-2k_2(\xi))$ is bounded on the compact $\bar B_1$: the coefficient map is piecewise linear, so the image of $\bar B_1$ lies in a real box $B'$, on which $\ln Z^a_{J_0}$ is convex (\cref{lem:s02:alg}(d)) and so bounded by its maximum over the vertices of $B'$, each finite by \Wtwo(a). For the second factor, $2\xi_+\in B''$ and $\ln Z^a_J(0,\cdot)$ is convex in the coefficients, so $Z^a_J(0,2k_1(\xi))\le\max_{v\in\vertex(B'')}Z^a_J(0,q_v)\le\max_v\sup_{J'\ge J_{q_v}}Z^a_{J'}(0,q_v)=: C_{B''}<\infty$ by \Wtwo\ at the $q_v$ (note $J\ge J_0\ge J_{q_v}$). By \cref{lem:s02:alg}(d), $|Z^a_J(\zeta)|\le Z^a_J(\Re\zeta)$, so $\sup_{J\ge J_0}\sup_{\cT_1}|Z^a_J|<\infty$. On $U\cap\cT_1$ the sources lie in $\cK$, so $J\mapsto Z^a_J(\zeta)$ is nondecreasing (first inequality of \cref{lem:s02:cut} at fixed $k$, then $k\to\infty$) and bounded, hence convergent. $U\cap\cT_1$ is a set of uniqueness for holomorphic functions on the connected $\cT_1$ (\cref{def:s08:enlarged}(iii)); Montel's theorem gives locally uniform convergence on $\cT_1$ to a holomorphic limit, which is $Z^a_\infty$ wherever the latter is defined (the same $J$-limit), and the monotone limit on $U\cap\cT_1$ gives $Z^a_J\le Z^a_\infty$ there. Under \Wtwo(b) for every $0\le q\in C^\infty_c$ the same holds on all of $\C^m$.
\end{proof}

\begin{lemma}[uniform local bounds for the modified-box functionals]\label{lem:s08:claim4}
Let $\cF$, $\Phi_k$, $\cT_1$, $T_\cF$ be as in \cref{def:s08:enlarged}.
\begin{itemize}[leftmargin=*,nosep]
\item[(A)] \textup{(Optional extension.)} Under \Wtwo(a) and \Wtwo(b) at the $q_v$, $v\in\vertex(B'')$: there is $C<\infty$, independent of $k$ and of $J\ge J_0$, with $\limsup_k\sup_{\cT_1}|\Phi_k|\le C$; and at fixed $J$, $\sup_k\sup_Q|\Phi_k|<\infty$ for every compact $Q\subset\C^m$ (this part uses \Wtwo(a) only).
\item[(A$_\infty$)] \textup{(Optional extension.)} Under \Wtwo\ for every $0\le q\in C^\infty_c$: the same with $\cT_1$ replaced by any compact $Q\subset\C^m$, with $C=C_Q$.
\item[(T)] \textup{(The tube.)} Under \cref{hyp:s03:W1}, \Itwo, \Ithree\ and the correlation inequalities of \cref{lem:s02:BFS}(a)--(e) only (no \Wtwo): the same for every compact $Q\subset T_\cF$ and all $J\ge J_0$, with no further threshold on $J$.
\item[(C)] \textup{(Curvature sector, $h\equiv0$; no hypothesis.)} $|\Phi_k(z)|\le\langle e^{F(f_{\Re z})}\rangle^{(3.38)}_{k,\Omega_{J,k}}\le\exp(\tfrac12\Var^{(3.38)}F(f_{\Re z}))$ and $\Var^{(3.38)}F(f)\le\sum_p\eps^3f_{c(p)}^2\to\norm f_{L^2}^2$, uniformly in $k$ and $J$.
\end{itemize}
\end{lemma}

\begin{proof}
$|\Phi_k(\zeta)|\le\Phi_k(\Re\zeta)$ by the proof of \cref{lem:s02:alg}(d) ($|\Theta(\zeta)|=\Theta(\Re\zeta)$ and the modulus of the multiplier), which is valid for any probability measure with the same real multiplier, in particular for the modified box state, which is not one of the states of \cref{def:s02:acube,def:s02:Kcube}. At a real $\xi$ the proof of \cref{lem:s02:alg}(f) (Cauchy--Schwarz with the multiplier identity, again valid for any such measure) gives
\begin{equation}\label{eq:s08:split}
\Phi_k(\xi)\ \le\ Z^{a,(3.38)}_{k,\Omega_{J,k}}(2f_\xi,-2k_2(\xi))^{1/2}\;Z^{a,(3.38)}_{k,\Omega_{J,k}}(0,2k_1(\xi))^{1/2}.
\end{equation}
\emph{First factor} (class $\cD$). By the $\cD$ half of \cref{lem:s07:L} with $C=C_J\cap\Lk$ (its hypothesis is \eqref{eq:s07:margin}(i)) it is $\le Z^a_{k,J}(2f_\xi,-2k_2)\le Z^a_{k,J_0}(2f_\xi,-2k_2)$, the second step being the second inequality of \cref{lem:s02:cut} at fixed $k$. The map $\xi\mapsto\eta(\xi)\defeq$ (coefficient vector of $(2f_\xi,-2k_2(\xi))$) is continuous, so $\eta(\Re Q)$, resp.\ $\eta(\bar B_1)$, lies in a real box $B'$; $\eta\mapsto\ln Z^a_{k,J_0}(\eta)$ is convex (\cref{lem:s02:alg}(d)), so $\sup_\xi Z^a_{k,J_0}(\eta(\xi))\le\max_{v\in\vertex(B')}Z^a_{k,J_0}(v)$, finitely many values, each convergent in $k$ by \cref{hyp:s03:W1} (the sources here have negative Wick coefficients, on which the theorem has no amplitude restriction; in versions (A), (A$_\infty$) by \Wtwo(a)), and convergence in $k$ implies boundedness in $k$. So the first factor is bounded uniformly in $k$, $\xi$ and $J\ge J_0$.
\emph{Second factor} (class $\cI$). By the $\cI$ half of \cref{lem:s07:U} with $B=C_{J+2n}\cap\Lk$ and $N\defeq\lceil R(\lmu)/\eps_k\rceil$ (hypothesis \eqref{eq:s07:margin}(ii)) it is $\le Z^a_{k,J+2n}(0,2k_1(\xi))$. (A): for $\xi\in B_1$, $2\xi_+\in B''$, so by convexity it is $\le\max_{v\in\vertex(B'')}Z^a_{k,J+2n}(0,q_v)$; each of these finitely many values converges in $k$ by \Wtwo(a) (hence is bounded in $k$, which is the fixed-$J$ bound Montel's theorem needs, and the $\limsup_k$ bound passes to any subsequential limit) with limit $Z^a_{J+2n}(0,q_v)\le\sup_{J'\ge J_{q_v}}Z^a_{J'}(0,q_v)<\infty$: this is \Wtwo\ at $q_v$, used directly ($J+2n\ge J_0\ge J_{q_v}$), uniformly in $J$; the finitely many instances $q_v$ are all of \Wtwo\ that version (A) uses. (A$_\infty$): the same with a box $B''\supset2(\Re Q)_+$ depending on $Q$. (T): for $\xi\in\Re Q$, $Q\subset T_\cF$, one has $0\le2k_1(\xi)\le\sigma_+/2\le s_1$ on the lattice (the sampled supremum is at most the continuum one), so \cref{cor:s11:W2prime}, $\ln N^a_{k,C_J}(0,q)\le\kappa\norm q_{1,k}$ for $0\le q\le s_1$ on $C_J\cap\Lk$, uniformly in $k$ and $J$, applied with $N^a_{k,X}(0,q)=Z^a_{k,X}(0,q)$ (the same functional with the same scheme multiplier, so no $J$-dependent factor intervenes), gives $Z^a_{k,J+2n}(0,2k_1(\xi))\le\exp(\kappa\norm{2k_1(\xi)}_{1,k})\le\exp\big(2\kappa\sum_j|\xi_j|\sup_{k'}\norm{\tilde h_j}_{1,k'}\big)$, bounded uniformly in $k$, $J$ and $\xi\in\Re Q$ because $\norm{\tilde h_j}_{1,k'}\to\norm{\tilde h_j}_{L^1}$ (Riemann sums, \eqref{eq:s02:riemann}). No lower bound on $J$ beyond $J_0$ enters: \cref{cor:s11:W2prime} holds for every $k$ and $J$.
\emph{(C).} At $h\equiv0$ the first inequality is \cref{lem:s02:alg}(d) as above and the second is \cref{lem:s02:GD} with $s=1$, the modified box state being a BFS measure (\cref{def:s07:box}). Its photon Gaussian has the unitary form for gauge-invariant observables (Step 1 of \cref{lem:s07:RPbox}), and $\Var^u\langle v,d_1A\rangle=\eps^{-1}\langle d_1^Tv,(d_1^Td_1+\lmu^2\eps^2)^{-1}d_1^Tv\rangle\le\eps^{-1}\norm v^2$; with $v_p=\eps^2f_{c(p)}$ this is $\Var F(f)\le\sum_p\eps^3f_{c(p)}^2$.
\end{proof}

\begin{remark}[one instance of \Wtwo\ in place of the vertex set]\label{rem:s08:oneinstance}
For $v\in\vertex(B'')$ one has $0\le q_v\le q_\zeta$ pointwise, and for $p=1+\eta>1$, $p'=p/(p-1)$, $q_v=\frac1p(pq_\zeta)-\frac1{p'}\big(p'(q_\zeta-q_v)\big)$. H\"older's inequality with the multiplier bookkeeping $e^{-a_k\sS(q_v)}=(e^{-a_k\sS(pq_\zeta)})^{1/p}(e^{a_k\sS(p'(q_\zeta-q_v))})^{1/p'}$ gives, for every state, $Z^a(0,q_v)\le Z^a(0,pq_\zeta)^{1/p}Z^a(0,-p'(q_\zeta-q_v))^{1/p'}$. The second factor is of class $\cD$, nonincreasing in $J$ (\cref{lem:s02:cut}, then $k\to\infty$ by \Wtwo(a)), hence $\le Z^a_{J_0}(0,-p'(q_\zeta-q_v))<\infty$, one real point of \cref{hyp:s03:W1}; the first is bounded uniformly in $J$ by \Wtwo(b) at the single source $(1+\eta)q_\zeta$. So \Wtwo\ at $(1+\eta)q_\zeta$, for any one $\eta>0$, implies the finitely many instances used in \cref{lem:s08:W1,lem:s08:claim4}(A). Only this direction is proved; no converse is claimed. Pointwise domination $q_v\le q_\zeta$ alone gives nothing, since $Z^a(0,\cdot)$ is not monotone in $q$: the Wick square has no sign, and the scheme multiplier $e^{-a_k\sS(q)}$ has no sign control.
\end{remark}

\subsection{King's state on the Wick tube}

\begin{proposition}[the Wick-tube object: existence and translations]\label{prop:s08:T}
Let $a$ be an admissible scheme (\cref{def:s03:adm}), and assume \Itwo\ and \Ithree. For a finite family $\cF'$ of real sources with nonnegative scalar members $\tilde h_j\ge0$ (any family is brought to this form by the splitting of \cref{def:s08:enlarged}) let $T_{\cF'}$ be its Wick tube, as in \cref{def:s08:enlarged}(i).
\begin{enumerate}[label=\textup{(\alph*)},leftmargin=*,nosep]
\item $Z^a_J\to Z^a_{\infty,T}$ locally uniformly on $T_{\cF'}$ as $J\to\infty$; $Z^a_{\infty,T}$ is holomorphic on $T_{\cF'}$; $Z^a_{\infty,T}=\Ginf$ on $T_{\cF'}\cap\Dtube$ and $Z^a_{\infty,T}=Z^a_\infty$ wherever the latter exists (the same $J$-limit of the same $Z^a_J$); $|Z^a_{\infty,T}(\zeta)|\le Z^a_{\infty,T}(\Re\zeta)$; $Z^a_{\infty,T}(\xi)>0$ for real $\xi\in T_{\cF'}$ with $-\xi\in T_{\cF'}$; and $Z^a_{\infty,T}(-z,w)=Z^a_{\infty,T}(z,w)$.
\item \textup{(Translations.)} $Z^a_{\infty,T}(\tau_b\cdot\zeta)=Z^a_{\infty,T}(\zeta)$ for all $b\in\R^3$ and $\zeta\in T_{\cF'}$, where $\tau_b$ acts on all sources; the translated family has the same tube.
\end{enumerate}
The proof uses, besides \cref{hyp:s03:W1}, only \cref{cor:s11:W2prime} (hence \Itwo, \Ithree\ and \cite[Thm.~6.2(1),(2)]{BFS1}), the second inequality of \cref{lem:s02:cut} (hence \cref{lem:s02:BFS}(d)), \cref{lem:s02:alg,lem:s02:baire}, Montel's theorem and the identity theorem; \Wtwo\ is not used.
\end{proposition}

\begin{proof}
(a) \emph{Local boundedness} of $\{Z^a_J\}_{J\ge J_0}$ on $T_{\cF'}$: $|Z^a_J(\zeta)|\le Z^a_J(\Re\zeta)\le Z^a_J(2f,-2k_2)^{1/2}Z^a_J(0,2k_1)^{1/2}\le Z^a_{J_0}(2f,-2k_2)^{1/2}\exp(\kappa\norm{k_1}_{L^1})$, by \cref{lem:s02:alg}(d),(f), the nonincreasing $\cD$-factor (second inequality of \cref{lem:s02:cut} at fixed $k$, then $k\to\infty$), and \cref{cor:s11:W2prime} in the limit $k\to\infty$ for the $\cI$-factor, valid since $2\norm{k_1}_\infty<\sigma_+/2\le s_1$; the first factor is bounded on compacts by the vertex argument of \cref{lem:s08:claim4} with \cref{hyp:s03:W1} at $C_{J_0}$. \emph{Convergence on $U\cap T_{\cF'}$}: there the values are real, nondecreasing in $J$ (first inequality of \cref{lem:s02:cut}, then $k\to\infty$) and bounded, since $Z^a_J(\mathrm if_\zeta,g_\zeta)\le Z^a_J(0,g_\zeta)$ by \cref{lem:s02:alg}(d) and $Z^a_J(0,g_\zeta)\le\exp(\kappa\norm{g_\zeta}_{L^1})$ by \cref{cor:s11:W2prime} in the limit $k\to\infty$ ($0\le g_\zeta\le\sigma_+/4\le s_1$ there). Montel's theorem on the connected open $T_{\cF'}$ with the uniqueness set $U\cap T_{\cF'}$ gives locally uniform convergence to a holomorphic $Z^a_{\infty,T}$; on $T_{\cF'}\cap\Dtube$ this is the limit $\Ginf$ of \cref{thm:s11:W}(i), the same argument on the other tube. The bound $|Z^a_{\infty,T}(\zeta)|\le Z^a_{\infty,T}(\Re\zeta)$ and the evenness in $z$ pass from the lattice (\cref{lem:s02:alg}(d); the invariance $(A,\vartheta)\mapsto(-A,-\vartheta)$ of every state (charge conjugation, \cref{def:s02:sources}), which maps $F(f)\mapsto-F(f)$ and fixes $r$; $T_{\cF'}$ is invariant under $z\mapsto-z$). Positivity at symmetric real points is \cref{lem:s02:alg}(e), $Z^a(\xi)Z^a(-\xi)\ge1$, in the limit.

(b) \emph{Lattice translations.} Let $b\in5^{-k_1}\Z^3$, $|b|_\infty\le M$, $k\ge k_1$. $\tau_b$ preserves $\Lk$ and $\sS_{\Lk}$, so $Z^a_{k,J}(\tau_b\cdot\zeta)=Z^a_{k,(C_J-b)\cap\Lk}(\zeta)$ (the action is translation covariant on $\Lk$). As sets of sites, $C_{J-1}\cap\Lk\subset(C_J-b)\cap\Lk\subset C_{J+1}\cap\Lk$, so the first inequality of \cref{lem:s02:cut} (stated for any finite site set) gives, for $\zeta\in U$ and $J\ge J_0+1$ (so that every plaquette and site at which the sources do not vanish lies in $C_{J-1}$), $Z^a_{k,J-1}(\zeta)\le Z^a_{k,J}(\tau_b\cdot\zeta)\le Z^a_{k,J+1}(\zeta)$. Let $k\to\infty$ and then $J\to\infty$ on $U\cap T_{\cF'}$, where the outer limits converge by (a): $Z^a_{\infty,T}(\tau_b\cdot\zeta)=Z^a_{\infty,T}(\zeta)$ there, hence on all of $T_{\cF'}$ by the identity theorem (both sides are holomorphic on the connected $T_{\cF'}$, the translated family having the same supremum norms). Iterating, the identity holds for all $b\in\bigcup_k5^{-k}\Z^3$.
\emph{All $b\in\R^3$.} For $\rho>0$ let $E_\rho\defeq C^\infty_c(\bar B_\rho;\Lambda^2)\times C^\infty_c(\bar B_\rho)$ (a Fr\'echet space) and $O_\rho\defeq\{(f,h)\in E_\rho:\ 4\norm h_\infty<\sigma_+/2\}$, open, convex and symmetric under $s\mapsto-s$. For $s=(f,h)\in O_\rho$ the one-member family with its splitting has the coefficient vector $(1,1,-1)$ in its tube ($2\norm{k_1}_\infty=2\norm{h^+}_\infty\le4\norm h_\infty<\sigma_+/2$), so $\Phi_T(s)\defeq\ln Z^a_{\infty,T}(f,h)$ is defined and finite on $O_\rho$ by (a). $\Phi_T$ is convex (a pointwise limit of the functions $\ln Z^a_{k,J}$, convex by H\"older's inequality) and Borel (a pointwise limit of continuous functions of finitely many sampled values). By \cref{lem:s02:baire}, $\Phi_T$ is continuous on $O_\rho$. For $s\in O_\rho$ and $b\in\R^3$ pick $b_n\to b$ in $\bigcup_k5^{-k}\Z^3$; then $\tau_{b_n}s,\tau_bs\in O_{\rho+|b|+1}$, $\tau_{b_n}s\to\tau_bs$ in $E_{\rho+|b|+1}$ (translations act continuously on $C^\infty_c$), and $\Phi_T(\tau_{b_n}s)=\Phi_T(s)$; so $\Phi_T(\tau_bs)=\Phi_T(s)$.
\emph{Complex coefficients.} For a family $\cF'$ the two holomorphic functions $\zeta\mapsto Z^a_{\infty,T}(\tau_b\cdot\zeta)$ and $\zeta\mapsto Z^a_{\infty,T}(\zeta)$ on $T_{\cF'}$ agree at the real points of $\{\xi\in\R^m:4\norm{h_\xi}_\infty<\sigma_+/2\}\cap T_{\cF'}$, where both limits exist and both equal the one-member value $\exp\Phi_T(f_\xi,h_\xi)$ (the same $J$-limit of the same $Z^a_J$). This intersection is open in $\R^m$ and contains $0$ (the first set alone need not lie in the tube: with $\tilde h_1=\tilde h_2$ and $\xi=(1,-1)$, $h_\xi=0$ while $k_1=\tilde h_1$), hence is a set of uniqueness in the connected $T_{\cF'}$; so the two functions agree on $T_{\cF'}$.
\end{proof}

\subsection{The two limits and the passage to every plane}

Throughout this subsection $\cF$, $\Phi_k$, $\zeta^*_{ij}$, $\cT_1$, $T_\cF$, $U$ are those of \cref{def:s08:enlarged} for a mid-plane $c\in D$, and ``version'' refers to \cref{lem:s08:claim4}: (C) is the curvature sector when the comparison family is convergent at fixed cubes (the uniform local bound being \cref{thm:s05:C}(b)), (T) the tube under \cref{hyp:s03:W1}, \Itwo, \Ithree, and (A), (A$_\infty$) the optional extension under \Wtwo.

\emph{Step 5 ($k\to\infty$ at fixed $J$).} By the fixed-$J$ bound of \cref{lem:s08:claim4} and Montel's theorem there is a subsequence $k_m\to\infty$ along which $\Phi_{k_m}\to\Phi^J$ locally uniformly on $\C^m$, with $\Phi^J$ entire. No convergence of the full sequence is claimed or needed. By \cref{lem:s07:RPbox}(b), $\sum_{ij}\bar\gamma_i\gamma_j\Phi_k(\zeta^*_{ij})\ge0$ for every $k\ge k_0(c,\delta)$ and $\gamma\in\C^N$; positive semidefiniteness of a Hermitian matrix is closed under entrywise limits, so $\sum_{ij}\bar\gamma_i\gamma_j\Phi^J(\zeta^*_{ij})\ge0$. The uniform bounds of \cref{lem:s08:claim4} pass to $\Phi^J$: $\sup_{\cT_1}|\Phi^J|\le C$ in (A), $\sup_Q|\Phi^J|\le C_Q$ in (A$_\infty$), (T) and (C), uniformly in $J\ge J_0$, whatever subsequence was taken. (The box $\Omega_{J,k}$ depends on $J$, so $\Phi^J$ does; the subsequence may depend on $J$.)

\emph{Step 6 ($J\to\infty$: sandwich on the cone, identity theorem off it).} By the $\cI$ halves of \cref{lem:s07:L,lem:s07:U}, whose hypotheses are \eqref{eq:s07:margin}(i),(ii), for every $k$ and every $\zeta\in U$,
\begin{equation}\label{eq:s08:sandwich}
Z^a_{k,J}(\zeta)\ \le\ \Phi_k(\zeta)\ \le\ Z^a_{k,J+2n}(\zeta).
\end{equation}
Let $k\to\infty$ along the subsequence of Step 5. The outer terms converge along the full sequence (version (C): \cref{thm:s05:C}(b); (T), (A), (A$_\infty$): \cref{hyp:s03:W1}, resp.\ \Wtwo(a), with Montel's theorem), so $Z^a_J\le\Phi^J\le Z^a_{J+2n}$ on $U$ for every subsequential $\Phi^J$. As $J\to\infty$ the outer terms converge on the respective uniqueness set: in (C) on $U=(\mathrm i\R)^{m_F}$ to $Z_\infty(\cdot,0)$ (\cref{thm:s05:C}(b)); in (T) on $U\cap T_\cF$ to $Z^a_{\infty,T}$ (\cref{prop:s08:T}(a); there $J\mapsto Z^a_J(\zeta)$ is real, nondecreasing and bounded by $Z^a_J(0,g_\zeta)\le\exp(\kappa\norm{g_\zeta}_{L^1})$, since $0\le g_\zeta\le\sigma_+/4\le s_1$); in (A) on $U\cap\cT_1$ to $Z^a_\infty$ (\cref{lem:s08:W1}); in (A$_\infty$) on $U$ to $Z^a_\infty$ (\cref{lem:s08:W1} with \Wtwo(b) for every $q$). Hence $\Phi^J\to Z^a_{\infty(,T)}$ pointwise on that set. By Step 5 the family $\{\Phi^J\}_J$ is locally bounded on the corresponding connected domain ($\C^{m}$, $T_\cF$, $\cT_1$, $\C^m$), on which the limit functional is holomorphic (\cref{thm:s05:C}(b), \cref{prop:s08:T}(a), \cref{lem:s08:W1}, \cref{lem:s08:W1} on $\C^m$ respectively); by Montel's theorem every subsequence of $(\Phi^J)$ has a further subsequence converging locally uniformly to a holomorphic function, which agrees with the limit functional on the uniqueness set and hence on the whole domain \cite{Hor90}. So $\Phi^J\to Z^a_{\infty(,T)}$ locally uniformly on the domain, in particular at the finitely many matrix points $\zeta^*_{ij}$ (which lie in $\cT_1$ by construction and in $T_\cF$ by \eqref{eq:s08:ampbound} when every $\amp(\zeta_i)<\sigma_W$), where $\Phi^J(\zeta^*_{ij})\to Z^a_{\infty(,T)}(\theta\zeta_i^*\sqcup\zeta_j)$ because the functionals depend on the family only through the total sources. Positivity passes to the limit:
\begin{equation}\label{eq:s08:midplane}
\sum_{i,j}\bar\gamma_i\gamma_j\,Z^a_{\infty(,T)}(\theta\zeta_i^*\sqcup\zeta_j)\ \ge\ 0\qquad\text{for the plane }\{x_\nu=c\},\ c\in D .
\end{equation}
A side result: the modified-box values converge to the cube state, $\lim_J\lim_{k_m}Z^{a,(3.38)}_{k_m,\Omega_{J,k_m}}(\zeta)=Z^a_{\infty(,T)}(\zeta)$, on the respective domain and along any $k$-subsequence on which the fixed-$J$ limits exist; this is the corrected form of King's passage from the modified boxes to the cube limit \cite[(3.42)]{K2}, with the margin $2n$ in place of one cube step (\cref{lem:s07:R,rem:s07:K8}). The reversal of King's second inequality is proved, weakly at every coupling and strictly at $e=0$ (\cref{lem:s07:R}); no step of the passage to the limit uses that inequality in King's direction.

\emph{Step 7 (every plane).} Let $c\in\R$ be arbitrary and the sources supported in $\{x_\nu>c\}$. If $c\in D$, \eqref{eq:s08:midplane} is the result. Otherwise pick $c'\in D$ and $b\defeq(c'-c)e_\nu$. The translated sources $\tau_b\zeta_i$ are supported in $\{x_\nu>c'\}$, and $\theta_{c'}(\tau_b\zeta)^*=\tau_b(\theta_c\zeta^*)$, because $\rho_{c'}(x+b)=\rho_c(x)+b$ for this $b$ and the linear parts of $\rho_c,\rho_{c'}$ agree. By \eqref{eq:s08:midplane} for the plane $c'$ and the translation invariance of the limit functional,
\[
\sum\bar\gamma_i\gamma_jZ^a_{\infty(,T)}(\theta_c\zeta_i^*\sqcup\zeta_j)=\sum\bar\gamma_i\gamma_jZ^a_{\infty(,T)}\big(\tau_b(\theta_c\zeta_i^*\sqcup\zeta_j)\big)=\sum\bar\gamma_i\gamma_jZ^a_{\infty(,T)}\big(\theta_{c'}(\tau_b\zeta_i)^*\sqcup\tau_b\zeta_j\big)\ \ge\ 0 .
\]
The translation invariance is supplied in (C) by \cref{prop:s06:translations}, and in (T) by \cref{prop:s08:T}(b) on $T_\cF$ (the amplitude $\amp$ is translation invariant, so the matrix points stay in the tube). In the optional extension a plane $c\notin D$ needs a non-lattice translation, which finitely many instances of \Wtwo(b) do not supply; there the step is carried by \cref{prop:s08:T}(b) together with \cref{lem:s08:W1} for $\cF$ and for $\tau_b\cF$ and the identity theorem on $\cT_1$ (\cref{thm:s23:OS2W}).

\subsection{The theorems}

\begin{theorem}[reflection positivity, curvature sector, comparison convention]\label{thm:s08:OS2C}
Assume that the comparison family is convergent at fixed cubes (\cref{def:s03:conv}). Let $\lmu>0$, $\nu\in\{1,2,3\}$, $c\in\R$, and let $\zeta_1,\dots,\zeta_N$ be finite families of real $2$-form sources $f_{il}\in C^\infty_c(\R^3;\Lambda^2)$ supported in the open half-space $\{x_\nu>c\}$, with arbitrary complex coefficients $z_{il}$ and no scalar sources. Then for all $\gamma\in\C^N$
\[
\sum_{i,j=1}^N\bar\gamma_i\gamma_j\,Z_\infty\big(\theta\zeta_i^*\sqcup\zeta_j,\,0\big)\ \ge\ 0 .
\]
Besides the convergence at fixed cubes the proof uses only \cite[Thm.~6.2(1),(2)]{BFS1} (through the first inequality of \cref{lem:s02:cut} and the $\cI$ halves of \cref{lem:s07:L,lem:s07:U}), \cref{thm:s05:C}(b) with \cref{lem:s05:C2}, and \cref{lem:s02:GD}, \cref{lem:s02:alg}(d),(e), \cref{lem:s02:BFS}(e), \cref{lem:s02:baire}; no inequality of class $\cD$ enters.
\end{theorem}

\begin{proof}
At $h\equiv0$ there is no Wick constant and no scheme, and $\cF$ consists of the $f_{il}$ and $\rho^*f_{il}$ with $U=(\mathrm i\R)^{m}$. What the argument of Steps 5--7 needs of the cube functionals is: the existence of $Z_J(\cdot,0)=\lim_kZ_{k,J}(\cdot,0)$ locally uniformly on $\C^m$, entire; the monotonicity of $J\mapsto Z_J(\mathrm if,0)$ on the cone (first inequality of \cref{lem:s02:cut}); the existence of $Z_\infty(\cdot,0)=\lim_JZ_J(\cdot,0)$, entire, locally uniformly on $\C^m$; and translation invariance. The first three are \cref{thm:s05:C}(b), whose local bound is the $\lmu$-free bound of \cref{thm:s05:C}(a) and \cref{lem:s05:C2} rather than a local bound resting on the nonincreasing clause of \cref{lem:s05:TWC} and so on \cite[Thm.~6.2(3),(4)]{BFS1}; the fourth is \cref{prop:s06:translations}, which uses \cref{lem:s02:baire} with $\ln Z_\infty$ finite on all of $E_\rho$ (upper bound from \cref{thm:s05:C}(b), lower bound \cref{lem:s02:alg}(e)) and the identity theorem. The lattice steps \cref{lem:s07:RPbox,lem:s07:L,lem:s07:U} include $g=0$, and only their $\cI$ halves are used; \cref{lem:s08:claim4}(C) is the fixed-$J$ and $J$-uniform bound; Steps 5--7 in version (C) give the claim.
\end{proof}

\begin{theorem}[reflection positivity on the Wick tube, comparison convention]\label{thm:s08:OS2WT}
Let $a$ be an admissible scheme, and assume \Itwo\ and \Ithree\ (the hypotheses of \cref{cor:s11:W2prime}); the $\cD$-class steps use in addition \cref{lem:s02:BFS}(d). Let $\sigma_+=s_1$ be the amplitude of \Itwo\ and $\sigma_W\defeq\sigma_+/8$. Let $\lmu>0$, $\nu$, $c\in\R$, and let $\zeta_1,\dots,\zeta_N$ be finite families of real sources $f_{il}\in C^\infty_c(\R^3;\Lambda^2)$, $h_{il}\in C^\infty_c(\R^3)$ supported in $\{x_\nu>c\}$, with complex coefficients $z_{il},w_{il}$ such that $\amp(\zeta_i)<\sigma_W$ for every $i$; in particular, all purely imaginary $w_{il}$ with arbitrary $z_{il}$ are admitted. Then the values $Z^a_{\infty,T}(\theta\zeta_i^*\sqcup\zeta_j)=\lim_JZ^a_J(\theta\zeta_i^*\sqcup\zeta_j)$ exist (\cref{prop:s08:T}) and
\[
\sum_{i,j=1}^N\bar\gamma_i\gamma_j\,Z^a_{\infty,T}\big(\theta\zeta_i^*\sqcup\zeta_j\big)\ \ge\ 0\qquad\text{for all }\gamma\in\C^N .
\]
Where $Z^a_\infty$ is defined, $Z^a_{\infty,T}=Z^a_\infty$ there. No lower bound on $J$ beyond $J_0$ enters, and \Wtwo\ is not used.
\end{theorem}

\begin{proof}
With the splitting of \cref{def:s08:enlarged} the matrix points lie in $T_\cF$ by \eqref{eq:s08:ampbound}. \Cref{lem:s08:claim4}(T) gives the bounds; Steps 5 and 6 in version (T), with \cref{prop:s08:T}(a) for the limit on $U\cap T_\cF$ and its holomorphy on $T_\cF$, give \eqref{eq:s08:midplane} for $c\in D$; Step 7 with \cref{prop:s08:T}(b) gives every plane. The $\cD$-class steps are the first factor of \eqref{eq:s08:split} (the $\cD$ half of \cref{lem:s07:L} and the second inequality of \cref{lem:s02:cut}) and the nonincreasing factor in \cref{prop:s08:T}(a).
\end{proof}

The arbitrary-amplitude form, reflection positivity of $Z^a_\infty$ at all complex coefficients under \Wtwo, is \cref{thm:s23:OS2W} in the optional extension; its proof is Steps 5--7 in version (A) for $c\in D$, where \Wtwo\ enters only at the finitely many $q_v$, and in version (A) or (A$'$) for $c\notin D$.

\begin{lemma}[reflection positivity of King's printed state]\label{lem:s08:OS2K}
\Cref{thm:s08:OS2C} holds for $\ZKinf$: for $\nu$, $c$, $\zeta_1,\dots,\zeta_N$ as there, $\sum_{i,j}\bar\gamma_i\gamma_j\ZKinf(\theta\zeta_i^*\sqcup\zeta_j,0)\ge0$ for all $\gamma\in\C^N$; and \cref{prop:s08:OS3,prop:s08:poly} below hold for $\ZKinf(\cdot,0)$. The inputs are \cref{hyp:s03:EC}, \cite[Thm.~6.2(1),(2)]{BFS1}, \cref{thm:s05:CK}(a),(b) in place of \cref{thm:s05:C}, \cref{lem:s04:D}(c) and \cref{rem:s04:F} for the two places where the cube convention is seen, and \cref{cor:s06:TK} for the translations; no inequality of class $\cD$ enters.
\end{lemma}

\begin{proof}
\Cref{lem:s07:RPbox,lem:s07:B} concern the modified box state alone and the lattice reflection, and \cref{lem:s08:claim4}(C) bounds $\Phi_k$ on the modified photon form: these are convention-free, as is Step 5 (Montel in $k$; positivity passes). The two places where the Neumann--Dirichlet cube states enter are the ends of the sandwich \eqref{eq:s08:sandwich}, \ie the $\cI$ halves of \cref{lem:s07:L,lem:s07:U}; on King's convention they read $Z^K_{k,J}(\zeta)\le\Phi_k(\zeta)\le Z^K_{k,J+2n}(\zeta)$ for every $k$ and $\zeta\in U=(\mathrm i\R)^m$: the lower end is \cref{lem:s07:L} with Step 1 of \cref{lem:s04:G} prepended (\cref{rem:s04:F}: $\avK{\cdot}\le\ava{\cdot}\le\langle\cdot\rangle^{(3.38)}_{k,\Omega_{J,k}}$ on $\cI$), the upper end is \cref{lem:s07:U} with $\avK{\cdot}$ on $X_{J+2n}$ in place of the comparison state (\cref{lem:s04:D}(b),(c): the box $\Omega_{J,k}$ does not meet the killed layers $T_{J+2n}$, and part (a) of \cref{lem:s07:U} is photon-only), under the hypotheses of \eqref{eq:s07:margin}. Here $J_0$ is the $J_0$ of the enlarged family (the sources and their reflections), which exceeds that of the sources alone when $c\ne0$; \cref{thm:s05:CK}(b) is applied to the enlarged family at $J$ and $J+2n$, $J\ge J_0$. Along Step 5's subsequence the ends converge along the full sequence by \cref{thm:s05:CK}(b), so $Z^K_J\le\Phi^J\le Z^K_{J+2n}$ on $U$ for every subsequential $\Phi^J$; as $J\to\infty$ both ends tend to $\ZKinf$ on $U$, $\{\Phi^J\}_J$ is locally bounded on $\C^m$ by \cref{lem:s08:claim4}(C), and $U$ is a set of uniqueness, so $\Phi^J\to\ZKinf$ locally uniformly on $\C^m$ as in Step 6, in particular at the matrix points; positivity passes to the limit. Step 7 uses \cref{cor:s06:TK} in place of \cref{prop:s06:translations}. \Cref{prop:s08:OS3,prop:s08:poly} use only the existence, holomorphy near $0$ and lattice symmetries of the limit, and the locally uniform convergences just established; their proofs below apply verbatim with $\ZKinf$, the charge-conjugation invariance being exact for $\avK{\cdot}$ as well.
\end{proof}

\subsection{\texorpdfstring{$\OSth$}{OS3} and the polynomial form}

\begin{proposition}[symmetry, realness and charge conjugation]\label{prop:s08:OS3}
Let $Z$ be one of: $Z_\infty(\cdot,0)$ if the comparison family is convergent at fixed cubes (resp.\ $\ZKinf(\cdot,0)$); $Z^a_{\infty,T}$ under the hypotheses of \cref{prop:s08:T}; $Z^a_\infty$ under \Wtwo\ (optional extension). The Schwinger functions
\[
S_{n,p}(f_1,\dots,f_n;h_1,\dots,h_p)\defeq\partial_{z_1}\cdots\partial_{z_n}\partial_{w_1}\cdots\partial_{w_p}\,Z\big(\textstyle\sum_iz_if_i,\sum_lw_lh_l\big)\big|_{0}
\]
(with $p=0$ and $S_n$ in the curvature sector) are (a) symmetric multilinear in the $f$'s and in the $h$'s; (b) real for real test functions; (c) zero for odd $n$. The proof uses only that $Z$ exists and is holomorphic near $0$, with the stated symmetries of the lattice functionals; the existence is the input of the respective sector.
\end{proposition}

\begin{proof}
(a) $Z$ depends on a family only through its total sources (so does every lattice functional, hence both limits), so its Taylor coefficients at $0$ are symmetric multilinear forms in the $f$'s and in the $h$'s; $Z^a_{\infty,T}$ is holomorphic on a tube containing $0$. The definition of $S_{n,p}$ does not depend on the family used to compute it (restrict a larger family to the given sources). (b) For real coefficients every lattice functional is real (real integrand and multiplier, \cref{lem:s02:alg}(a)), so $Z$ is real at the real points near $0$ (for $Z^a_{\infty,T}$ on $\R^m\cap T$, an open piece of $\R^m$ containing $0$); a holomorphic function real on the real points near $0$ has real Taylor coefficients at $0$. (c) Every cube state and every modified box state is invariant under $(A,\vartheta)\mapsto(-A,-\vartheta)$ (the photon terms are even in $A$, the masked divergence included; $\cos(\vartheta_x-e\eps A-\vartheta_y)$ is invariant; $W^a$ depends on $r$ only; Lebesgue measure is preserved), which maps $F_{\Lk}(f)\mapsto-F_{\Lk}(f)$ and fixes $W^a_{\Lk}(h)$. Hence $Z_{k,X}(-z,w)=Z_{k,X}(z,w)$ for all $k$, $X$, and in both limits $Z(-z,w)=Z(z,w)$ (on the $z\mapsto-z$ invariant tube for $Z^a_{\infty,T}$). So $\partial_z^{(n)}\partial_w^{(p)}Z(0)=0$ for odd $n$. This is the argument of \cite[lem:C-OS3lat, thm:C-OS3]{SymPhase} for the cube states.
\end{proof}

\begin{proposition}[the polynomial form of reflection positivity]\label{prop:s08:poly}
Let $P=\sum_\alpha c_\alpha\prod_iF(f_{\alpha i})\prod_lW^a(h_{\alpha l})$ be a polynomial in the smeared fields with real sources supported in $\{x_\nu>c\}$, and $\theta P\defeq\sum_\alpha\bar c_\alpha\prod_iF(\rho^*f_{\alpha i})\prod_lW^a(h_{\alpha l}\circ\rho)$. Then $S(\theta(P)P)\ge0$, where $S$ assigns to a monomial the corresponding Schwinger function of \cref{prop:s08:OS3}: the Taylor coefficient at $0$ of $Z^a_{\infty,T}$ under the hypotheses of \cref{thm:s08:OS2WT}, resp.\ of $Z_\infty(\cdot,0)$ (of $\ZKinf(\cdot,0)$), where at $h\equiv0$ the $W^a$ factors are absent.
\end{proposition}

\begin{proof}
Let all sources of $P$ be supported in $\{x_\nu\ge c+\delta\}$ and let $\cF'$ be the family of all sources of $P$ and their reflections. Its scalar members have no sign, so Steps 4--6 are applied to the split family $\cF'_{\rm split}$ (members $h^\pm_{\alpha l}$, $h^\pm_{\alpha l}\circ\rho$, all $\ge0$); its tube contains $0$ ($k_1(0)=0$), and the unsplit functional is the composition of the split one with the linear map $w_{\alpha l}\mapsto(w^+_{\alpha l},w^-_{\alpha l})=(w_{\alpha l},-w_{\alpha l})$, so by the chain rule the derivatives at $0$ of the unsplit $\Phi_k$, $\Phi^J$, $Z^a_{\infty(,T)}$ are finite linear combinations of those of the split ones, and locally uniform convergence near $0$ passes through the linear map. (i) \emph{Lattice.} For $k\ge k_0$ the sampled $P$ is a gauge-invariant polynomial in the plus-side smeared fields, lies in $L^2$ of the modified box state (Gaussian moments; $\lambda>0$) and in $\mathfrak B_+$ (\cref{def:s07:reflection}), and $\theta(P)=\overline{P\circ\hat\rho}$ by \eqref{eq:s07:thetaTheta}. By \cref{lem:s07:RPbox}(a) (its $L^2$ extension), $\langle\theta(P)P\rangle^{(3.38)}_{k,\Omega_{J,k}}\ge0$. (ii) \emph{Moments as derivatives.} $\langle\theta(P)P\rangle^{(3.38)}=\sum_{\alpha,\beta}\bar c_\alpha c_\beta\langle\prod_iF(\rho^*f_{\alpha i})\prod_lW^a(h_{\alpha l}\circ\rho)\prod_iF(f_{\beta i})\prod_lW^a(h_{\beta l})\rangle^{(3.38)}$, and each mixed moment is a mixed partial derivative at $0$ of $\Phi_k$ for the family $\cF'$ (differentiate under the integral; the scheme multiplier $e^{-wa_k\sS(h)}$ is the $-a_k$ inside $W^a$). (iii) \emph{Limits.} By Steps 5--6 for $\cF'_{\rm split}$, $\Phi_{k_m}\to\Phi^J$ and $\Phi^J\to Z^a_{\infty(,T)}$ locally uniformly on a neighbourhood of $0$, so by Weierstrass' theorem all derivatives at $0$ converge, and $\sum\bar c_\alpha c_\beta\,\partial^{(\alpha,\beta)}Z^a_{\infty(,T)}(0)\ge0$; the left side is $S(\theta(P)P)$ by the family-independence of the Schwinger functions (\cref{prop:s08:OS3}(a)). Every plane: Step 7 (translation of the sources; $S$ is translation invariant by \cref{prop:s08:T}(b), resp.\ \cref{prop:s06:translations} or \cref{cor:s06:TK}). The curvature sector is the case $h\equiv0$ with no splitting. This is the form OS2(b) of \cite[def:OS, thm:C-OS2]{SymPhase} for the cube states.
\end{proof}

\begin{remark}[which form is used, and the amplitude constants]\label{rem:s08:amplitudes}
The main theorem (\cref{thm:s13:main}) states reflection positivity of the Wick sector in the tube form of \cref{thm:s08:OS2WT}, with \cref{prop:s08:OS3,prop:s08:poly} for the Schwinger functions; its inputs for this clause are \cref{hyp:s03:W1}, \Itwo, \Ithree\ and \cite[Thm.~6.2(1)--(4)]{BFS1}, without \Wtwo. This suffices for the consumer: the weak form of $\OSf$ and the reconstruction of \cref{thm:s09:S2K} use reflection positivity for exponentials $e^{\mathrm i\Phi(f)}$ of the smeared fields, whose Wick coefficients are purely imaginary, so amplitude $0$ only. Three positive, $\eps$-free amplitude constants occur in the Wick sector: $\sigma_W=\sigma_+/8$ here; $H/2=\sigma_+/32$, the negative-part radius of the tube $\Dtube$ of \cref{thm:s11:W} ($H=s_1/16$); and $\sigma_0=H/4$ in the bound $\Ezz$ of \cref{thm:s11:W}(ii). On $T_\cF\cap\Dtube$ the objects $Z^a_{\infty,T}$ and $\Ginf$ coincide (\cref{prop:s08:T}(a)). On King's convention the tube statements of this section are read under \TWK\ (\cref{rem:s11:Kconvention}, \cref{prop:s13:WK}); the curvature-sector statement needs no such hypothesis (\cref{lem:s08:OS2K}). The margin $R(\lmu)$ of \cref{lem:s07:U} enters the proof, not the conclusion: reflection positivity is proved at every fixed $\lmu>0$, and in a limit of vanishing photon mass it passes to the limit functionals as a pointwise limit of positive semidefinite matrices, the polynomial form by Weierstrass' theorem.
\end{remark}
